\begin{center}
\textsc{Gauss} to \textsc{Olbers}.\quad Göttingen, 14~January 1827.
\end{center}

\par\noindent\hspace*{\parindent}\mbox{.\ .\ .\ .\ .\ .}\quad I believe I told you before that I have found it quite impracticable to make the work I intend to publish on my measurements entirely self-contained in theoretical respects as well, unless I first treat at least a large part of the theory separately elsewhere.
Explaining the grounds for my operations will already make for a bulky work, but, if I may put it this way, the \textls[120]{grounds} for the \textls[120]{grounds} cannot go into the work itself without making it quite a patchwork and still leaving it unsatisfactory.
I have therefore decided first to develop various theoretical subjects in separate papers, since only in this way can I work out these important new chapters of mathematics with a certain degree of completeness.
In a sense my prize essay on the transformation of surfaces is the first of these papers. I submitted the second to the royal Society a few months ago as Supplementum theoriae combinationis observationum etc.
It contains the principles that must be applied as the basis for the \textls[120]{adjustment} of the observations, and even a few examples from \textsc{Krayenhoff}'s measurements and my own.\newline
\mbox{.\ .\ .\ .\ .\ .}\par
